\documentclass[11pt,a4paper]{article}
\usepackage[margin=2.2cm]{geometry}
\usepackage[T1]{fontenc}
\usepackage[utf8]{inputenc}
\usepackage{booktabs}
\usepackage{longtable}
\usepackage{graphicx}
\usepackage{amsmath}
\usepackage{hyperref}
\usepackage{setspace}
\usepackage{microtype}
\usepackage{array}
\usepackage{xcolor}

\setstretch{1.08}
\setlength{\parskip}{4pt}
\setlength{\parindent}{0pt}

\hypersetup{
  colorlinks=true,
  linkcolor=blue,
  urlcolor=blue,
  pdftitle={KOA Risk Prediction: Internal Validation and Preliminary Cross-Cultural Replication},
  pdfauthor={BioAge-KOA-Prediction Team}
}

\title{Predicting Symptomatic Knee Osteoarthritis Risk from CHARLS Data:\\
\large Internal Validation with Early Unseen Holdout, Calibration, and Decision Analysis,\\
with Preliminary Cross-Cultural Replication on the US Health and Retirement Study}
\author{BioAge-KOA-Prediction Team}
\date{April 2026}

\begin{document}
\maketitle

\begin{abstract}
Symptomatic knee osteoarthritis (KOA) is a major source of disability and health-care burden in aging
populations. We developed and internally validated a machine-learning prediction model for symptomatic KOA
in a CHARLS-derived cohort ($n=12{,}329$, KOA prevalence $13.3\%$) and subsequently applied it to the US
Health and Retirement Study (HRS, RAND Longitudinal File 2022, Wave 13) as a preliminary cross-cultural
replication exercise.

\textbf{Development and internal validation.}\ Three design choices underpinned the CHARLS pipeline:
(i) an early 80/20 stratified unseen holdout split before model development,
(ii) train-fitted missing-data imputation to avoid leakage, and
(iii) a pre-specified primary model (Random Forest, $n_{\text{trees}}=300$) applied to 17 raw
sociodemographic and clinical variables (raw17). Reporting follows TRIPOD guidance and the PROBAST
framework. The primary model achieved ROC-AUC $0.847$ / PR-AUC $0.688$ / Brier $0.070$ / ECE $0.045$
on the internal out-of-fold estimate, and ROC-AUC $0.907$ / PR-AUC $0.791$ / Brier $0.053$ / ECE $0.056$
on the unseen holdout. Decision-curve analysis (DCA) showed strictly positive net benefit in the
10\%--30\% threshold band relative to both treat-all and treat-none strategies
(${\approx}64$ avoided unnecessary interventions per 100 patients on the holdout).

\textbf{Preliminary HRS replication.}\ The trained model was applied to $n=20{,}605$ HRS Wave 13
respondents aged $\geq 45$, using 12 of the 16 raw17 predictors that could be harmonised from the
publicly available RAND fat file (raw16; Biological Age, Residence, Dyslipidemia, and wave identifier
were unavailable). The KOA proxy was defined as current physician-diagnosed arthritis plus difficulty
stooping/kneeling (prevalence $33.9\%$). ROC-AUC fell to $0.604$, calibration slope to $0.184$, and DCA
net benefit was negative relative to treat-all. This substantial attenuation is attributable to three
compounding factors: the exclusion of Biological Age (the dominant raw17 predictor), the use of a
broader surrogate outcome definition, and genuine cross-cultural and demographic differences between
Chinese and US older adult populations. These results do not constitute formal external validation; they
frame the minimum expected performance degradation and motivate the full harmonisation plan
(ELSA as primary external validation cohort, with raw17 including KDM-computed Biological Age).
\end{abstract}

\textbf{Keywords:} symptomatic KOA, machine learning, TRIPOD, calibration, decision curve analysis,
SHAP, internal validation, external validation, HRS, cross-cultural replication

% ============================================================
\section{Introduction}
% ============================================================
Knee osteoarthritis is one of the leading causes of years lived with disability worldwide, with
age-standardised prevalence having risen steadily in recent decades\cite{Safiri2020}. In aging
populations, symptomatic KOA -- defined as clinically diagnosed knee osteoarthritis co-occurring with
knee pain -- imposes substantial mobility, quality-of-life, and health-care-utilization costs.
Risk-prediction models may support earlier stratification, but publication-quality evidence requires
more than high discrimination alone. Reporting should include calibration behaviour and threshold-level
decision consequences, not only ROC-AUC, and should follow standards such as TRIPOD\cite{Collins2015}
and PROBAST\cite{Wolff2019}.

A recent CHARLS-based report\cite{Fu2025} presented strong discrimination for symptomatic KOA and
highlighted biological age as a central predictor. However, translating a published result to a new
analytical pipeline is sensitive to cohort filtering, missing-data strategy, categorical encoding,
class-balance handling, and split design. This manuscript frames the project as a single
internal-validation narrative with a pre-specified primary model and explicit internal-evidence
boundaries, and extends the analysis with a preliminary cross-cultural replication on the RAND Health
and Retirement Study (HRS) Longitudinal File to characterise minimum transportability expectations.

% ============================================================
\section{Methods}
% ============================================================

\subsection{Study Objective and Outcome}
The prediction target was symptomatic KOA in a CHARLS-derived cohort. The outcome was operationalised
from two CHARLS-derived fields as
\[
\text{KOA} = \mathbb{1}\bigl[\text{Arthritis} = \text{``yes''}\bigr]\ \wedge\
\mathbb{1}\bigl[\text{position\_knees} = \text{``yes''}\bigr],
\]
\noindent which matches the ``symptomatic KOA'' operationalisation used by Fu et al.\cite{Fu2025}
(physician-reported arthritis with knee-localised joint pain). Both source fields were removed from the
predictor set to prevent target leakage; rows with \texttt{position\_knees} missing were dropped because
the target itself could not be verified for those subjects.

\subsection{Study Cohort and Exclusion Cascade}
The analysis was built on Sheet1 of the CHARLS extract (15{,}545 baseline rows, 28 columns). Applying
the target-verification rule removed 3{,}216 rows with \texttt{position\_knees} missing, yielding the
analytical cohort ``Dataset A'' with 12{,}329 rows and 13.32\% symptomatic-KOA prevalence. Dataset A
was then partitioned once, before any model development, into an internal training partition
(80\%, $n=9{,}863$, prevalence 13.32\%) and an unseen holdout partition (20\%, $n=2{,}466$, prevalence
13.30\%) using a stratified split with \texttt{random\_state=42}. Table~\ref{tab:table1} summarises
the characteristics of Dataset A and its two partitions.

\begin{table}[h!]
\centering
\small
\caption{Cohort characteristics of the CHARLS-derived Dataset A and the early 80/20 partitions.
Values are counts (\%) or mean $\pm$ SD.}
\label{tab:table1}
\begin{tabular}{lccc}
\toprule
Characteristic & Full Dataset A & Internal train (80\%) & Unseen holdout (20\%) \\
\midrule
N & 12{,}329 & 9{,}863 & 2{,}466 \\
Symptomatic KOA, n (\%) & 1{,}642 (13.32) & 1{,}314 (13.32) & 328 (13.30) \\
Biological Age, mean $\pm$ SD & $58.70 \pm 9.79$ & $58.70 \pm 9.76$ & $58.70 \pm 9.90$ \\
Age bucket $\geq 60$, n (\%) & 3{,}357 (27.2) & 2{,}676 (27.1) & 681 (27.6) \\
Female, n (\%) & 6{,}443 (52.3) & 5{,}143 (52.1) & 1{,}300 (52.7) \\
BMI, mean $\pm$ SD & $23.65 \pm 3.84$ & $23.63 \pm 3.83$ & $23.73 \pm 3.88$ \\
BMI overweight, n (\%) & 10{,}122 (82.1) & 8{,}081 (81.9) & 2{,}041 (82.8) \\
BMI obese, n (\%) & 1{,}488 (12.1) & 1{,}192 (12.1) & 296 (12.0) \\
Married, n (\%) & 10{,}429 (84.6) & 8{,}343 (84.6) & 2{,}086 (84.6) \\
Education low/medium/high, \% & 28.1 / 61.9 / 10.0 & 27.9 / 62.0 / 10.1 & 28.8 / 61.4 / 9.7 \\
Residence urban, n (\%) & 7{,}954 (64.5) & 6{,}388 (64.8) & 1{,}566 (63.5) \\
Hypertension, n (\%) & 3{,}142 (25.5) & 2{,}503 (25.4) & 639 (25.9) \\
Dyslipidemia, n (\%) & 1{,}108 (9.0) & 882 (8.9) & 226 (9.2) \\
Diabetes, n (\%) & 726 (5.9) & 584 (5.9) & 142 (5.8) \\
Cancer, n (\%) & 93 (0.8) & 75 (0.8) & 18 (0.7) \\
CVD, n (\%) & 1{,}304 (10.6) & 1{,}016 (10.3) & 288 (11.7) \\
Smoking ever, n (\%) & 4{,}896 (39.7) & 3{,}897 (39.5) & 999 (40.5) \\
Drinking ever, n (\%) & 5{,}144 (41.7) & 4{,}156 (42.1) & 988 (40.1) \\
\bottomrule
\end{tabular}
\end{table}

\subsection{Feature Set (raw17)}
The pre-specified primary feature set, referred to throughout as \emph{raw17}, contains 17 variables
(Table~\ref{tab:raw17}). Secondary feature branches extending raw17 with PhenoAge clocks or broader
KDM bundles are retained as sensitivity analyses only.

\begin{table}[h!]
\centering
\small
\caption{The 17 variables in the raw17 feature set used by the primary model.}
\label{tab:raw17}
\begin{tabular}{llp{6.2cm}}
\toprule
Variable & Type & CHARLS-derived meaning \\
\midrule
wave & integer & Survey wave identifier \\
Time & integer & Survey year indicator \\
Gender & binary & Self-reported sex (1=male, 2=female) \\
Age\_New & ordinal (2) & Age bucket (45--59, $\geq 60$) \\
Marital & binary & Marital status (married vs other) \\
Education & ordinal (3) & Educational attainment (low/medium/high) \\
Residence & binary & Urban vs rural \\
Hypertension & binary & Self-reported physician-diagnosed \\
Dyslipidemia & binary & Self-reported physician-diagnosed \\
Diabetes & binary & Self-reported physician-diagnosed \\
Cancer & binary & Self-reported physician-diagnosed \\
CVD & binary & Self-reported cardiovascular disease \\
Smoke & binary & Ever vs never smoker \\
Drink & binary & Ever vs never drinker \\
BMI & continuous & Measured body-mass index (kg/m\textsuperscript{2}) \\
BMI\_New & ordinal (3) & BMI category (underweight/normal, overweight, obese) \\
Biological Age & continuous & CHARLS-supplied biological-age field \\
\bottomrule
\end{tabular}
\end{table}

\subsection{Missing Data Handling}
After the target-verification exclusion, the raw17 feature set and the outcome were complete for every
row in Dataset A (0\% missingness). As a pipeline safety net and to support sensitivity branches with
additional biomarker features, missing values were handled with train-fitted simple imputation (median
for numeric, most-frequent for categorical). Imputers and standardisation (\texttt{StandardScaler})
were fit on training data only and then applied to validation or unseen data.

\subsection{Early Unseen Holdout Strategy}
The early 80/20 stratified split was created before model development and persisted as CSV artefacts
(\texttt{dataset\_A\_internal\_train.csv}, \texttt{dataset\_A\_unseen\_test.csv}). Internal model
development and OOF estimation used only the internal partition; the unseen partition was reserved for
a single-shot external-style evaluation.

\subsection{Primary and Secondary Modeling Strategy}
The pre-specified primary model was a Random Forest (\texttt{RandomForestClassifier},
\texttt{n\_estimators=300}, \texttt{random\_state=42}, all other hyperparameters at scikit-learn
defaults). No grid search, Bayesian search, or Optuna-based tuning was performed on this configuration.
KDM-based and adapted aging-clock feature branches are retained as secondary sensitivity analyses.

\subsection{Class-Imbalance Handling}
No resampling (SMOTE, random under-/over-sampling) and no \texttt{class\_weight} adjustment were
applied to the primary Random Forest. PR-AUC was pre-specified as the primary discrimination metric
alongside ROC-AUC, following standard recommendations for imbalanced binary prediction.

\subsection{Evaluation Framework}
Internal validation used stratified 5-fold out-of-fold (OOF) estimation on the internal training
partition, followed by re-evaluation on the unseen holdout. Reported metrics:
\begin{itemize}
\item discrimination: ROC-AUC, PR-AUC;
\item calibration: Brier score, ECE, MCE, Cox logistic calibration slope and intercept,
      decile reliability diagram;
\item decision relevance: Vickers net benefit at thresholds 0.05--0.60 and threshold-band
      summaries for 0.10--0.30.
\end{itemize}
Probability calibration was assessed for both uncalibrated (``raw'') and isotonic-calibrated
(\texttt{CalibratedClassifierCV}, \texttt{method=``isotonic''}, 3-fold inner CV) variants.

\subsection{Interpretability}
Model interpretation included SHAP global importance\cite{Lundberg2017}, SHAP interaction analysis,
and local waterfall plots for representative cases.

\subsection{Differences vs.\ Reference Study (Fu et al., 2025)}
Table~\ref{tab:paper-delta} summarises the key methodological differences between the reference
study and the present pipeline.

\begin{table}[h!]
\centering
\caption{Methodological delta versus reference study.}
\label{tab:paper-delta}
\begin{tabular}{p{3.4cm}p{5.1cm}p{5.1cm}}
\toprule
Domain & Reference paper & Current pipeline \\
\midrule
Cohort framing & Symptomatic KOA in CHARLS-derived sample &
  Symptomatic KOA in CHARLS-derived sample; updated internal protocol with early holdout \\
Missing data & Study-specific exclusions / imputation choices &
  Train-fitted \texttt{SimpleImputer} within each split \\
Split protocol & 70/30 emphasised & Early 80/20 unseen holdout + internal OOF validation \\
Class-balance handling & Model-specific handling &
  No resampling, no \texttt{class\_weight} adjustment \\
Primary metric & ROC-AUC-centred & PR-AUC reported alongside ROC-AUC \\
Primary model & XGBoost-centred & Random Forest + raw17 pre-specified \\
\bottomrule
\end{tabular}
\end{table}

% ============================================================
\subsection{Preliminary HRS Cross-Cultural Replication}
% ============================================================

\subsubsection*{HRS Cohort and Data Source}
The Health and Retirement Study (HRS)\cite{HRS2023} is a US nationally representative longitudinal
panel of adults aged $\geq 50$ funded by the National Institute on Aging. For this replication we used
the RAND HRS Longitudinal File 2022 (Version 1)\cite{RAND2023}, a derived analysis-ready dataset
spanning survey waves 1992--2022. Wave 13 (survey year 2016) was selected as the target wave,
corresponding to the most recent wave with broad variable coverage in the RAND fat file.

\subsubsection*{HRS Eligibility and KOA Proxy Definition}
Inclusion criteria: respondents with Wave 13 interview data and age $\geq 45$ years
(\texttt{r13iwstat} indicating completed interview; $n=20{,}912$ before age filter,
$n=20{,}640$ after). Because the RAND fat file does not contain a direct symptomatic KOA indicator,
a two-component proxy was constructed:
\begin{enumerate}
  \item \textbf{Current physician-reported arthritis:} \texttt{r13arthr} = 1 (``reports arthritis
        this wave''), capturing active disease rather than lifetime ever-diagnosis;
  \item \textbf{Knee-relevant functional impairment:} \texttt{r13stoopa} = 1 (``any difficulty
        stooping, kneeling, or crouching''), chosen as the closest available HRS analogue to
        knee-localised pain.
\end{enumerate}
\[
\text{KOA\_HRS} = \mathbb{1}\bigl[\texttt{r13arthr}=1\bigr]\ \wedge\
\mathbb{1}\bigl[\texttt{r13stoopa}=1\bigr]
\]
This proxy results in a prevalence of 33.9\% ($n_{\text{pos}}=6{,}976$), substantially higher than
the CHARLS KOA prevalence of 13.3\%. The difference reflects the broader specificity of the HRS proxy
relative to the CHARLS knee-localised pain criterion, the inclusion of inflammatory arthritis types
captured by \texttt{r13arthr}, and genuine cross-national prevalence differences between Chinese and
US elderly populations. Rows with missing arthritis response were excluded ($n=35$), yielding a
final analytic sample of $n=20{,}605$.

\subsubsection*{HRS Feature Harmonisation (raw16)}
Of the 17 raw17 features, 12 could be directly harmonised from the RAND fat file
(Table~\ref{tab:hrs-harmonisation}). Four variables could not be mapped and were excluded from this
replication: (1) \emph{Biological Age} is not provided in the RAND fat file (requires linkage to the
HRS Venous Blood Study biomarker file and KDM-BA computation); (2) \emph{Residence} (urban/rural)
requires the restricted-access HRS geographic crosswalk and was set to a single constant (urban=1),
introducing non-differential noise; (3) \emph{Dyslipidemia} (\texttt{r13cholst}) was absent from
Wave 13 of the RAND fat file; and (4) \emph{wave identifier} is not applicable across cohorts.
The resulting 12-variable feature set is labelled \emph{raw16} throughout. The trained CHARLS model
was re-fitted on the raw16 subset of the CHARLS internal training partition and then applied to HRS.

\begin{table}[h!]
\centering
\small
\caption{Feature harmonisation map between CHARLS raw17 and HRS RAND fat-file variables for Wave 13.
Entries marked --- indicate variables not available in the RAND public file; entries marked [const]
were set to a fixed constant as a conservative fallback.}
\label{tab:hrs-harmonisation}
\begin{tabular}{lllp{4cm}}
\toprule
raw17 variable & CHARLS source & HRS variable & Notes \\
\midrule
Gender & Derived & \texttt{ragender} & 1=male, 2=female (same coding) \\
Age\_New & Derived & \texttt{r13agey\_b} & 1 if 45--59, 2 if $\geq 60$ \\
Marital & Derived & \texttt{r13mstat} & 1=married, 2=other \\
Education & Derived & \texttt{raedyrs} & years $\rightarrow$ 3-level ordinal \\
Residence & Derived & --- [const=1] & Restricted file needed \\
Hypertension & Self-report & \texttt{r13hibpe} & 1=ever diagnosed \\
Dyslipidemia & Self-report & --- & Not in Wave 13 RAND file \\
Diabetes & Self-report & \texttt{r13diabe} & 1=ever diagnosed \\
Cancer & Self-report & \texttt{r13cancre} & 1=ever diagnosed \\
CVD & Self-report & \texttt{r13hearte} & 1=ever heart disease \\
Smoke & Self-report & \texttt{r13smokev} & 1=ever smoked regularly \\
Drink & Self-report & \texttt{r13drinkd} & 1=drinks \\
BMI & Measured & \texttt{r13bmi} & kg/m\textsuperscript{2}, same units \\
BMI\_New & Derived & Derived from \texttt{r13bmi} & Same 3-level categorisation \\
Biological Age & Computed & --- & Requires HRS VBS linkage \\
wave / Time & Admin & --- & Not applicable cross-cohort \\
\bottomrule
\end{tabular}
\end{table}

\subsubsection*{HRS Statistical Analysis}
The same evaluation pipeline applied to the CHARLS holdout was applied to HRS: ROC-AUC, PR-AUC,
Brier score, ECE, MCE, Cox logistic calibration slope and intercept, a decile reliability diagram,
and Vickers DCA in the 10\%--30\% threshold band. No recalibration was applied before the primary
assessment; observed miscalibration is reported as a transparency artefact.

% ============================================================
\section{Results}
% ============================================================

\subsection{Primary Model Performance (Random Forest + raw17)}
Table~\ref{tab:primary-metrics} summarises the matched internal OOF and unseen-holdout performance
for the pre-specified primary model. Discrimination remained strong in both settings, and
unseen-holdout ROC-AUC and PR-AUC were numerically higher than internal OOF estimates. Brier scores
and ECE values were low; the Cox calibration slope and intercept indicate residual miscalibration to
be interpreted in an internal-validation context.

\begin{table}[h!]
\centering
\small
\caption{Primary model (Random Forest + raw17) metrics. Slope/intercept are from Cox logistic
recalibration; see text for the isotonic caveat.}
\label{tab:primary-metrics}
\begin{tabular}{llccccccc}
\toprule
Set & Variant & ROC-AUC & PR-AUC & Brier & ECE & Slope & Intercept & F1 \\
\midrule
Internal OOF & Raw & 0.847 & 0.688 & 0.070 & 0.045 & 0.989 & $-0.040$ & 0.649 \\
Internal OOF & Isotonic & 0.839 & 0.676 & 0.069 & 0.042 & 1.607 & 1.000 & 0.653 \\
Unseen holdout & Raw & 0.907 & 0.791 & 0.053 & 0.056 & 1.216 & 0.178 & 0.801 \\
Unseen holdout & Isotonic & 0.900 & 0.779 & 0.051 & 0.060 & 1.742 & 1.088 & 0.799 \\
\bottomrule
\end{tabular}
\end{table}

\subsection{Reliability Diagnostics and the Isotonic-Slope Anomaly}
The Cox calibration-slope values for the isotonic-calibrated variants ($1.607$ internal, $1.742$
holdout) appear paradoxical against the expectation that isotonic recalibration should move slopes
toward 1.0. However, the slope/intercept reported by logistic-recalibration-in-the-large is defined
on the continuous logit of predicted probabilities and is known to be ill-behaved when applied to
piecewise-constant isotonic output\cite{NiculescuMizil2005,VanCalster2019}. ECE actually
\emph{decreases} after isotonic calibration in the internal OOF setting ($0.045 \rightarrow 0.042$).
The decile reliability table (Table~\ref{tab:reliability-oof}) and the reliability diagrams
(supplementary artefacts) are the appropriate evidence for the isotonic variant.

\begin{table}[h!]
\centering
\small
\caption{Decile reliability of the primary Random Forest model on the internal OOF set
($n=9{,}863$).}
\label{tab:reliability-oof}
\begin{tabular}{lrrrrr}
\toprule
Bin (raw) & n & Mean predicted & Observed rate & $|$obs $-$ pred$|$ \\
\midrule
$(-0.001, 0.007]$ & 1{,}311 & 0.003 & 0.028 & 0.025 \\
$(0.007, 0.013]$  &   719   & 0.012 & 0.033 & 0.022 \\
$(0.013, 0.027]$  & 1{,}101 & 0.021 & 0.036 & 0.015 \\
$(0.027, 0.040]$  &   827   & 0.035 & 0.041 & 0.006 \\
$(0.040, 0.063]$  & 1{,}101 & 0.053 & 0.045 & 0.007 \\
$(0.063, 0.090]$  &   921   & 0.077 & 0.062 & 0.016 \\
$(0.090, 0.130]$  &   964   & 0.110 & 0.082 & 0.028 \\
$(0.130, 0.197]$  &   961   & 0.162 & 0.107 & 0.054 \\
$(0.197, 0.383]$  &   971   & 0.267 & 0.162 & 0.105 \\
$(0.383, 0.930]$  &   987   & 0.629 & 0.743 & 0.114 \\
\bottomrule
\end{tabular}
\end{table}

\subsection{Decision-Curve Analysis}
Decision-curve analysis follows the net-benefit framework of Vickers and
colleagues\cite{Vickers2006,Vickers2016}:
\[
\mathrm{NB}_{\text{model}}(t) = \frac{\mathrm{TP}}{n} -
\frac{\mathrm{FP}}{n} \cdot \frac{t}{1-t}.
\]
Across the 10\%--30\% threshold band, the primary model produced strictly positive mean net benefit
against both treat-all and treat-none. The derived net-reduction-per-100-patients
figure\cite{Vickers2016} was $56.0$ (internal OOF) and $63.8$ (unseen holdout) avoided unnecessary
interventions versus treat-all; for isotonic probabilities the corresponding values were $56.5$ and
$65.8$, with relative reductions of $65.2\%$ and $75.9\%$ against the treat-all baseline.

\subsection{Secondary Sensitivity Branches}
Secondary feature branches did not exceed the primary raw17 profile: \texttt{raw17\_plus\_pheno}
achieved ROC-AUC $0.827$ (internal) / $0.880$ (unseen); \texttt{raw17\_plus\_adapted\_aging}
achieved $0.741$ / $0.775$.

% ============================================================
\subsection{Preliminary HRS Replication Results}
% ============================================================

Table~\ref{tab:hrs-metrics} presents the performance of the CHARLS-trained model (re-fitted on raw16)
applied to the HRS Wave 13 analytic sample ($n=20{,}605$).

\begin{table}[h!]
\centering
\small
\caption{Performance of the CHARLS-trained model (Random Forest + raw16) applied to HRS Wave 13
($n=20{,}605$, KOA proxy prevalence 33.9\%) compared with CHARLS internal results. The HRS
analysis uses 12 of 17 raw17 features; Biological Age, Residence, Dyslipidemia, and wave
identifier are excluded. Calibration slope outside [0.80, 1.20] indicates systematic miscalibration;
logistic recalibration-in-the-large is recommended before any clinical application in this
population.}
\label{tab:hrs-metrics}
\begin{tabular}{lccccccc}
\toprule
Setting & $n$ & KOA prev & ROC-AUC & PR-AUC & Brier & ECE & Cal slope \\
\midrule
CHARLS -- Internal OOF  & 9{,}863  & 13.3\% & 0.847 & 0.688 & 0.070 & 0.045 & 0.989 \\
CHARLS -- Unseen holdout & 2{,}466  & 13.3\% & 0.907 & 0.791 & 0.053 & 0.056 & 1.216 \\
\midrule
HRS Wave 13 (raw16) & 20{,}605 & 33.9\% & 0.604 & 0.399 & 0.258 & 0.194 & 0.184 \\
\bottomrule
\end{tabular}
\end{table}

Discrimination fell substantially from the CHARLS holdout to HRS
(ROC-AUC: $0.907 \rightarrow 0.604$; PR-AUC: $0.791 \rightarrow 0.399$). Calibration was severely
degraded (slope $0.184$, well outside the acceptable $[0.80,\,1.20]$ range), reflecting the
prevalence mismatch between CHARLS (13.3\%) and the HRS KOA proxy (33.9\%): the model, calibrated
for a 13\% base rate, systematically under-predicts risk in a population where the true (proxy) event
rate is 34\%. DCA net benefit versus treat-all was negative across the 10\%--30\% threshold band
(mean gain $-0.072$, corresponding to $-33.8$ avoided interventions per 100 patients), confirming
that the uncalibrated model offers no decision-support advantage over treating all patients in this
setting.

\paragraph{Interpretation.} Three factors jointly explain the observed performance gap:

\begin{enumerate}
  \item \textbf{Missing Biological Age.} Biological Age is the dominant predictor in the CHARLS
        raw17 model (SHAP analysis, Step 04). Its exclusion from raw16 removes the strongest
        discriminative signal and likely accounts for the majority of the ROC-AUC loss.
  \item \textbf{Outcome-definition heterogeneity.} The HRS KOA proxy (current arthritis plus
        difficulty stooping/kneeling) is broader than the CHARLS criterion (arthritis plus
        knee-localised pain). Outcome noise inflates apparent false-positive rates and degrades
        calibration independently of feature quality.
  \item \textbf{Cross-cultural and demographic differences.} The HRS sample is predominantly
        non-Hispanic White US adults; CHARLS is a Chinese nationally-representative cohort.
        Differences in health behaviours, diagnostic ascertainment culture, comorbidity
        distributions, and healthcare access are expected to attenuate cross-cultural
        transportability.
\end{enumerate}

These HRS results do not constitute formal external validation; they represent a conservative lower
bound on transportability under maximum harmonisation constraints and explicitly motivate the full
validation plan described in the Discussion.

% ============================================================
\section{Discussion}
% ============================================================
This manuscript frames symptomatic KOA prediction as a methodologically transparent
internal-validation study, with a preliminary cross-cultural replication on HRS as a transportability
stress test. Primary claims are anchored to a single pre-specified model-feature combination (Random
Forest + raw17), reporting follows TRIPOD\cite{Collins2015} with a completed checklist supplement,
and internal risk-of-bias was assessed against PROBAST\cite{Wolff2019}.

\paragraph{Internal validation.}
The primary raw17 model achieved ROC-AUC 0.847 (OOF) / 0.907 (holdout) and PR-AUC 0.688 / 0.791,
substantially exceeding the prevalence baseline. The reliability diagnostics
(Table~\ref{tab:reliability-oof}) identify systematic under-prediction in the top decile of predicted
risk -- a pattern attributable to the bounded-probability property of tree ensembles\cite{NiculescuMizil2005}
-- and strictly positive DCA net benefit in the 10\%--30\% threshold band provides a grounded
internal-evidence basis for risk-stratification decision support.

\paragraph{Comparison with the reference study.}
The reference study by Fu et al.\cite{Fu2025} reported XGBoost-centred discrimination with biological
age as the dominant predictor, using a 70/30 split. The present pipeline replicates the same outcome
operationalisation and confirms that the raw17 feature set supports strong discrimination even without
explicit biomarker enrichment. The methodological delta (Table~\ref{tab:paper-delta}) shows that the
main structural difference is the adoption of an early stratified holdout and OOF validation rather
than a single held-out test, which better separates development decisions from final performance claims.

\paragraph{HRS replication and transportability.}
The HRS replication (ROC-AUC 0.604; calibration slope 0.184) reveals substantial cross-cultural
performance attenuation and must be interpreted carefully. The three factors identified in the Results
-- missing Biological Age, outcome heterogeneity, and cross-cultural differences -- are hierarchically
ordered: Biological Age alone is expected to account for the majority of the observed ROC-AUC gap
based on its SHAP importance rank. Two downstream consequences follow. First, the HRS result
establishes a \emph{lower bound} on cross-cultural transportability under worst-case harmonisation;
actual performance with full raw17 (including KDM-computed Biological Age from HRS Venous Blood Study
biomarker data) is expected to be substantially higher. Second, the severe miscalibration (slope
$0.184$) is mathematically predictable from the prevalence mismatch: a model calibrated to $13\%$
applied to a $34\%$ prevalence population will systematically produce predictions centred near
$13\%$ when the correct centre is $34\%$; logistic recalibration-in-the-large (re-estimating the
intercept on a small HRS anchor sample) should restore acceptable calibration\cite{Steyerberg2019}.

\paragraph{External validation and generalisability.}
The most critical gap separating this study from submission-strength clinical claims remains the
absence of a true full external validation. We define a two-stage plan: (A) \emph{HRS full raw17
replication} -- obtain the HRS Venous Blood Study biomarker file, compute KDM-based Biological Age
via \texttt{step\_03\_kdm\_ba.py}, and re-run the replication with the complete raw17 feature set
alongside logistic recalibration-in-the-large; and (B) \emph{ELSA primary external validation} --
apply the model to the English Longitudinal Study of Ageing (ELSA), which includes direct biomarker
assays enabling full raw17 harmonisation and a knee-pain module enabling a more specific KOA outcome
proxy. The Korean Longitudinal Study of Ageing (KLoSA) and WHO SAGE cohorts are retained as
secondary targets. Until Stage B is completed, clinical-utility language remains explicitly limited
to internal-validation decision-support signals.

\paragraph{Strengths.}
Strengths include the large population-based cohort, transparent leakage controls, pre-specification
of the primary model before development, multi-dimensional performance evaluation (discrimination,
calibration, clinical utility, interpretability), adherence to TRIPOD and PROBAST, and the addition
of a cross-cultural replication that explicitly characterises the model's transportability envelope.

% ============================================================
\section{Limitations}
% ============================================================
\begin{itemize}

\item \textbf{Internal evidence only (primary).} The CHARLS internal validation uses a single
national cohort; performance on geographically, ethnically, or institutionally distinct populations
is formally unknown.

\item \textbf{Unseen holdout is not a true external cohort.} Participants in both partitions are
drawn from the same data source and share the same confounder distribution; the holdout estimates a
lower bound on expected performance degradation rather than true generalisability.

\item \textbf{HRS replication is preliminary and feature-incomplete.} The HRS analysis uses raw16
(12 of 17 features), excludes the dominant predictor (Biological Age), uses a surrogate KOA
outcome, and applies the model to a cross-cultural population without recalibration. The HRS result
should not be interpreted as a formal estimate of model performance in US populations.

\item \textbf{Outcome-definition heterogeneity across cohorts.} CHARLS KOA (arthritis $+$ knee pain)
differs from the HRS proxy (arthritis $+$ stooping difficulty). Outcome noise degrades apparent
performance independently of the predictive model.

\item \textbf{Self-reported comorbidities.} Hypertension, diabetes, cancer, CVD, dyslipidemia, and
lifestyle variables are self-reported in CHARLS, inheriting recall bias and diagnostic ascertainment
heterogeneity.

\item \textbf{Bootstrap confidence intervals not yet tabulated.} Point estimates and reliability
diagrams are the current evidence basis for calibration; bootstrap CIs for discrimination and
calibration metrics will be added in the next revision (TRIPOD item 16).

\item \textbf{Cross-sectional outcome.} The KOA label is cross-sectionally derived and cannot
distinguish incident from prevalent cases in the CHARLS extract.

\item \textbf{Isotonic calibration slope anomaly.} Cox calibration-slope estimates for the
isotonic-calibrated variant should be interpreted alongside ECE and the decile reliability diagram,
not in isolation, as explained in the reliability-diagnostics section.

\end{itemize}

% ============================================================
\section{Conclusion}
% ============================================================
This study demonstrates a methodologically rigorous internal-validation pipeline for symptomatic KOA
prediction in a large Chinese nationally-representative cohort. The pre-specified primary model
(Random Forest + raw17) achieves strong discrimination (ROC-AUC 0.907 on the unseen holdout),
acceptable calibration, and positive decision-curve net benefit in the clinically relevant
10\%--30\% risk threshold band, meeting TRIPOD reporting standards throughout.

A preliminary cross-cultural replication on the US Health and Retirement Study (HRS, Wave 13,
$n=20{,}605$) reveals substantial performance attenuation under maximum harmonisation constraints
(ROC-AUC 0.604, calibration slope 0.184). This attenuation is largely explained by the unavailability
of Biological Age in the RAND public file, the use of a surrogate KOA proxy, and genuine
cross-cultural differences, rather than by a fundamental failure of the raw17 feature set. The
findings establish a conservative lower bound on model transportability and motivate a two-stage
full external validation plan: HRS full raw17 replication (Stage A) and ELSA primary external
validation (Stage B). Until Stage B is completed, clinical-utility claims remain bounded to the
internal CHARLS evidence.

% ============================================================
\section*{Data Availability}
% ============================================================
The CHARLS development dataset (``Raw Data of Biological Age'' derived from CHARLS) is publicly
available under a CC BY 4.0 licence at Mendeley Data:
\url{https://data.mendeley.com/datasets/3rv7mf5pv9/1} (Fu F, Mendeley Data V1, 2025;
doi:10.17632/3rv7mf5pv9.1). The RAND HRS Longitudinal File is publicly available at
\url{https://hrsdata.isr.umich.edu/data-products/rand} upon free registration. The complete
analysis pipeline code is available at the project repository.

% ============================================================
\begin{thebibliography}{99}
% ============================================================

\bibitem{Fu2025} Fu F, Dong L, Lian J, et al. Biological age threshold is associated with
symptomatic knee osteoarthritis risk in Chinese adults: insights from machine learning analysis
of a national cohort. \textit{PLOS One}. 2025;20(12):e0335250.

\bibitem{Collins2015} Collins GS, Reitsma JB, Altman DG, Moons KGM. Transparent reporting of a
multivariable prediction model for individual prognosis or diagnosis (TRIPOD): the TRIPOD statement.
\textit{BMJ}. 2015;350:g7594.

\bibitem{Wolff2019} Wolff RF, Moons KGM, Riley RD, et al. PROBAST: a tool to assess the risk of
bias and applicability of prediction model studies. \textit{Ann Intern Med}. 2019;170(1):51--58.

\bibitem{Zhao2014} Zhao Y, Hu Y, Smith JP, Strauss J, Yang G. Cohort profile: the China Health and
Retirement Longitudinal Study (CHARLS). \textit{Int J Epidemiol}. 2014;43(1):61--68.

\bibitem{Safiri2020} Safiri S, Kolahi AA, Smith E, et al. Global, regional and national burden of
osteoarthritis 1990--2017: a systematic analysis of the Global Burden of Disease Study 2017.
\textit{Ann Rheum Dis}. 2020;79(6):819--828.

\bibitem{Vickers2006} Vickers AJ, Elkin EB. Decision curve analysis: a novel method for evaluating
prediction models. \textit{Med Decis Making}. 2006;26(6):565--574.

\bibitem{Vickers2016} Vickers AJ, van Calster B, Steyerberg EW. Net benefit approaches to the
evaluation of prediction models, molecular markers, and diagnostic tests. \textit{BMJ}. 2016;352:i6.

\bibitem{VanCalster2019} Van Calster B, McLernon DJ, van Smeden M, Wynants L, Steyerberg EW.
Calibration: the Achilles heel of predictive analytics. \textit{BMC Med}. 2019;17(1):230.

\bibitem{NiculescuMizil2005} Niculescu-Mizil A, Caruana R. Predicting good probabilities with
supervised learning. In: \textit{Proceedings of the 22nd International Conference on Machine
Learning (ICML)}. 2005:625--632.

\bibitem{Steyerberg2019} Steyerberg EW. \textit{Clinical Prediction Models: A Practical Approach
to Development, Validation, and Updating}. 2nd ed. Springer; 2019.

\bibitem{Lundberg2017} Lundberg SM, Lee SI. A unified approach to interpreting model predictions.
In: \textit{Advances in Neural Information Processing Systems}. 2017.

\bibitem{Brier1950} Brier GW. Verification of forecasts expressed in terms of probability.
\textit{Monthly Weather Review}. 1950;78(1):1--3.

\bibitem{Breiman2001} Breiman L. Random forests. \textit{Machine Learning}. 2001;45:5--32.

\bibitem{Chen2016} Chen T, Guestrin C. XGBoost: a scalable tree boosting system. In:
\textit{Proceedings of the 22nd ACM SIGKDD International Conference on Knowledge Discovery
and Data Mining}. 2016:785--794.

\bibitem{Klemera2006} Klemera P, Doubal S. A new approach to the concept and computation of
biological age. \textit{Mechanisms of Ageing and Development}. 2006;127(3):240--248.

\bibitem{Zadrozny2002} Zadrozny B, Elkan C. Transforming classifier scores into accurate multiclass
probability estimates. In: \textit{Proceedings of the 8th ACM SIGKDD International Conference on
Knowledge Discovery and Data Mining}. 2002.

\bibitem{HRS2023} Health and Retirement Study. \textit{Public Use Dataset}. Produced and distributed
by the University of Michigan with funding from the National Institute on Aging (grant number
NIA U01AG009740). Ann Arbor, MI; 2023.

\bibitem{RAND2023} RAND HRS Longitudinal File 2022 (V1). Santa Monica, CA: RAND Corporation; 2023.
\url{https://hrsdata.isr.umich.edu/data-products/rand}.

\end{thebibliography}

\end{document}
